\chapter{Objetivos}

En este capítulo se presentan el objetivo general y los objetivos específicos que orientaron el desarrollo de la investigación. El grado de cumplimiento de cada uno de ellos se evalúa de forma explícita en el capítulo de conclusiones.

\section{Objetivo general}

Desarrollar y validar un marco metodológico para el análisis, la implementación y la evaluación comparativa de algoritmos metaheurísticos, clásicos e híbridos, aplicados al dimensionamiento óptimo de sistemas híbridos de energías renovables que integran generación eólica, generación fotovoltaica, almacenamiento electroquímico y producción de hidrógeno verde, sustentando las conclusiones sobre desempeño algorítmico en evidencia experimental replicada y auditable.

\section{Objetivos específicos}

\begin{enumerate}

\item \textbf{Analizar} la estructura matemática de un modelo representativo de sistema híbrido de energías renovables con producción de hidrógeno, identificando sus variables de decisión, funciones objetivo, restricciones operacionales y modelos físicos de conversión energética.

\item \textbf{Implementar y verificar} computacionalmente dicho modelo mediante pruebas de conservación y límites, y cuantificar su grado de reproducción respecto de indicadores publicados mediante errores absolutos y relativos.

\item \textbf{Evaluar críticamente} los supuestos y simplificaciones presentes en el modelo de referencia, identificando sus limitaciones en la representación de la variabilidad del recurso renovable, del comportamiento del almacenamiento y de las restricciones de interacción con la red.

\item \textbf{Formular} un conjunto de instancias de optimización con un espacio de diseño ampliado respecto del modelo de referencia, definiendo su codificación, su función objetivo, su estrategia de tratamiento de restricciones y el óptimo exacto de cada instancia sobre el dominio discretizado.

% PENDIENTE-AVAL: redacción modificada, requiere aval del profesor guía
\item \textbf{Implementar y comparar} el desempeño de un conjunto de metaheurísticas de población clásicas y de un híbrido construido sobre ellas (un enjambre de partículas con operadores genéticos, PSO+GA) sobre dichas instancias, bajo un presupuesto computacional idéntico y frente a una línea base no informada, analizando la calidad de las soluciones, el comportamiento de convergencia y la consistencia entre ejecuciones independientes.

% PENDIENTE-AVAL: redacción modificada, requiere aval del profesor guía
\item \textbf{Caracterizar} el desempeño relativo de los algoritmos mediante estadística descriptiva de sus ejecuciones independientes (mediana, dispersión, brecha respecto del óptimo exacto y tasa de éxito) y un protocolo de comparación estadística fijado a priori, con dos familias de contrastes: el híbrido frente a cada uno de los demás algoritmos, en particular frente a los dos métodos de los que parte y frente a la línea base sin aprendizaje, y cada metaheurística clásica frente a esa línea base; tanto dentro de cada instancia como sobre el conjunto de instancias.

\item \textbf{Determinar} la sensibilidad de los resultados a las restricciones operacionales, al tratamiento del año bisiesto y a la extrapolación del recurso eólico, explicitando el alcance de cada comparación.

\end{enumerate}
